\section{Purpose, conventions and status}

\paper{} denotes P.~Enflo, \emph{On the invariant subspace problem in Hilbert spaces},
arXiv:2305.15442v2 (source in \texttt{paper/}). Equation numbers ``\paper{} (17)'' refer to it.

This document has two jobs. It restates, with complete proofs, the parts of \paper{} that are
correct or can be corrected (Sections \ref{sec:reductions}--\ref{sec:variation}). And it tracks
the parts that are not proofs (Section~\ref{sec:gaps}), together with evidence about whether they
can be repaired (Sections~\ref{sec:model}--\ref{sec:experiments}). It does not claim the theorem.
For the focused statement of the continuation obstruction, start with the standalone paper
\texttt{tex/refutation.tex}; these notes retain the broader investigation and its open questions.

\subsection*{Status tags}
\begin{itemize}[leftmargin=2em]
\item \stproved: a complete proof is written in this document.
\item \stoutline: the statement is believed and the argument is sketched, but details or
constants have not been written out.
\item \stopen: no proof is known to us.
\item \stlean{name}: additionally machine-checked in the Lean project at the repository root.
\end{itemize}

\subsection*{Conventions}
$H$ is a complex Hilbert space and $\ip{\cdot}{\cdot}$ is linear in the \emph{first} argument, as in
\paper{}. (Mathlib's inner product is linear in the second, so $\ip{a}{b}$ here is
\texttt{inner b a} in Lean.) $\ell^2=\ell^2(\N_0)$ with unilateral shift $S$, identified with the
Hardy space $H^2$ via $a\mapsto a(z)=\sum_j a_jz^j$. We write $\kappa=\norm{T}$ and assume
$0<\kappa\le\frac12$ where \paper{} takes $\kappa=10^{-20}$; nothing so far needs more.
We write $\tau$ for the quantity \paper{} calls $\varepsilon\theta$.

\subsection*{Summary of where things stand}
\begin{center}
\small
\begin{tabularx}{\textwidth}{@{}>{\raggedright\arraybackslash}p{.40\textwidth}
>{\raggedright\arraybackslash}p{.12\textwidth}>{\raggedright\arraybackslash}X@{}}
\toprule
Part of \paper{} & Here & Status\\
\midrule
Reductions (p.~1) & \S\ref{sec:reductions} & proved\\
Extremal series, (1)--(18) & \S\ref{sec:extremal} & proved; two lemmas in Lean\\
Type 1 / type 2, choice of $x_0$, Lemma 1 & \S\ref{sec:start} & proved after correction\\
Lemmas 2--5 & \S\ref{sec:domination} & proved after correction\\
Case I, (25) & \S\ref{sec:scaling} & proved\\
Case II, (26)--(27) & \S\ref{sec:scaling} & (27) fails under the stated small-defect hypotheses; defect decrease open\\
Variation formulas (28)--(30) & \S\ref{sec:variation} & proved after correction\\
Main Construction (31)--(45) & \S\ref{sec:gaps}, Gaps \ref{gap:secondorder}--\ref{gap:real} & open; experiments in \S\ref{sec:experiments}\\
Convergence step, (46) and after & \S\ref{sec:rigidity}, Gap \ref{gap:indep} & argument refuted in a model (Theorem~\ref{thm:nogo})\\
Type 2: reduction, (20), (47) & \S\ref{sec:type2} & proved / outline\\
Type 2: endgame & \S\ref{sec:type2}, Gap \ref{gap:type2} & open; depends on the type 1 construction\\
\bottomrule
\end{tabularx}
\end{center}
